% META: {"id":"1SPE-GEOREP-EX-047","chapitre":"1SPE-GEOMETRIE-REPEREE","type_objet":"exercice","capacites":["1SPE-GEOMETRIE-REPEREE-C5"],"capacites_codes":["C5"],"methodes":["M5"],"parcours":3,"competences":["calculer","raisonner"],"duree_min":20,"mode_creation":"ex_nihilo","sources_inspiration":[],"fichier_tex":"chapitres/1SPE-GEOMETRIE-REPEREE/exercices/1SPE-GEOREP-EX-047.tex","corrige_tex":"chapitres/1SPE-GEOMETRIE-REPEREE/corriges/1SPE-GEOREP-CO-047.tex","status":"approved","gestes":["raisonnement_par_coordonnees"],"origin":"generated","status_history":["generated","audited","accepted","approved"],"release_acceptance":"RELEASE_OWNER_BATCH_ACCEPTANCE","acceptance_closure_digest":"sha256:9b3ccf9a81c5520fbb7e03b7b2d3e7bf2057a02834b6d908bf3be5f49f3b553f","acceptance_receipt_digest":"sha256:4fad86f0282e0fa4e0419cca1ad6379ae7af5a280c04a4a90880f90a1c044242"}
% BEGIN-VERIFY
% from sympy import *
% x, y = symbols('x y')
% # Orthocentre du triangle A(0,0), B(6,0), C(2,4)
% # Hauteur issue de A: perpendiculaire a BC. BC=(−4,4), normal a la hauteur
% # -4x + 4y = 0 => x - y = 0
% # Hauteur issue de B: perpendiculaire a AC. AC=(2,4), normal a la hauteur
% # 2(x-6) + 4(y-0) = 0 => 2x + 4y - 12 = 0 => x + 2y - 6 = 0
% sol = solve([x-y, x+2*y-6], [x,y])
% assert sol == {x: 2, y: 2}
% END-VERIFY
\begin{exercice}{1SPE-GEOREP-EX-047}{3}{20}
Soit le triangle $ABC$ avec $A(0 ; 0)$, $B(6 ; 0)$ et $C(2 ; 4)$. Déterminer les coordonnées de l'orthocentre $H$.
\end{exercice}
